\clearpage
\setcounter{table}{0}
\renewcommand{\thetable}{A\arabic{table}}
\setcounter{figure}{0}
\renewcommand{\thefigure}{A\arabic{figure}}

\section*{Appendix}
\addcontentsline{toc}{section}{Appendix}

\section{Data Construction Details}\label{app:data}

\paragraph{Source and scope.} The CoreLogic Owner Transfer file compiles
recorded deeds from county recorder offices in all fifty states and the
District of Columbia. Our extract contains all record types---deeds of sale,
quitclaims, intra-family conveyances, trust transfers, and foreclosure
instruments. Coverage density is stable from 2007 onward; recording lags
truncate the file around August 2024. All annual statistics use 2007--2023;
the Proposition 19 cohort analyses use outcome windows that close by June 30,
2024.

\paragraph{Deduplication.} Multi-parcel transactions and re-recorded
corrections generate multiple rows per economic event. We keep one record per
parcel-date pair, preferring arm's-length category records and, within
category, the highest recorded consideration.

\paragraph{Classification.} The taxonomy in Section~\ref{sec:data} is
implemented with the recorder-derived interfamily indicator, the
arm's-length category code, deed type, recorded consideration, and
string features of the buyer and seller name fields (surname equality;
exact same-person matches on either buyer slot, which define the
same-person retitle class; trust keywords; estate, executor, probate, and
administrator keywords).
Name fields are populated for the large majority of records (the seller
surname is non-missing for 84 percent of interfamily events and the buyer
surname for 66 percent); missing names can only push true family transfers
into the ``other interfamily'' class, never into the market class, so the
same-surname series is a strict lower bound on inter-person family transfers.

\paragraph{Absentee recipients.} A recipient is absentee when the buyer
mailing ZIP on the deed differs from the property's situs ZIP, among records
where both are observed (77 percent of events in 2007, rising to 98 percent
by 2023). Within-state comparisons are unaffected by the secular improvement
in coverage; difference-in-differences designs absorb it in time effects.

\section{Robustness Checks}\label{app:robustness}

\paragraph{Trust-inclusive volume DiD.} Re-estimating the Proposition 19
volume difference-in-differences with trust self-transfers included in the
family aggregate yields a CA $\times$ Post coefficient of $-0.241$
(s.e.\ $0.027$), a 21 percent decline, confirming that the steady-state drop
in family transfers is not an artifact of reclassification between direct
family deeds and trust vehicles.

\paragraph{Permutation inference.} Because California is the only treated
state, the text reports placebo-state permutation $p$-values: each donor
state is assigned the treatment in turn, each design is re-estimated, and
California's coefficient is ranked in the absolute-value distribution across
all 51 estimates. For the family-volume DiD, California ranks 5th of 51
($p \approx 0.10$); for the placebo-netted estimate, 9th ($p \approx 0.18$);
for the sold-within-24-months triple difference, $p = 0.67$; for the
estate/trust-seller market-sale DiD, $p = 0.71$ (levels) and $p = 0.63$
(net of overall market).

\paragraph{Raw difference-in-differences means.} Table~\ref{tab:raw22}
reports the unadjusted means underlying the sold-within-24-months estimates.

\begin{table}[!ht]\centering
\caption{Sold on open market within 24 months: raw cohort means}
\label{tab:raw22}
{\small
\begin{tabular}{llrrr}
\toprule
Cohort group & Region & Pre (2017m1--2018m12) & Post (2021m7--2022m6) & Change \\
\midrule
Family transfer & California & 16.1\% \; (n=470{,}163) & 12.3\% \; (n=247{,}417) & $-3.8$pp \\
Family transfer & Rest of US & 13.1\% \; (n=1{,}793{,}673) & 11.1\% \; (n=1{,}118{,}028) & $-2.0$pp \\
Market purchase & California & 6.2\% \; (n=788{,}627) & 4.7\% \; (n=421{,}807) & $-1.5$pp \\
Market purchase & Rest of US & 6.9\% \; (n=7{,}722{,}553) & 6.5\% \; (n=4{,}495{,}875) & $-0.4$pp \\
\bottomrule
\end{tabular}}
\begin{minipage}{0.95\textwidth}\vspace{2pt}{\scriptsize \textit{Notes:}
Non-corporate recipients only. Family transfer = same-surname,
cross-surname interfamily, and estate/executor classes (same-person retitles
excluded). The difference-in-differences for family transfers is $-1.8$pp;
for market purchases, $-1.1$pp; the triple difference is $-0.74$pp
(regression counterpart in Table~\ref{tab:prop19}; permutation $p = 0.67$).}\end{minipage}
\end{table}
